Las métricas usadas para evaluar la clasificación se presentan a continuación, donde los valores a calcular dependen de los verdaderos positivos $TP$, verdaderos negativos $TN$, falsos positivos $FP$, y falsos negativos $FN$, asociados a la clase de cada ejemplo del conjunto de datos.


\begin{itemize}
    \item Exactitud: esta medida corresponde a la proporción de predicciones predichas correctamente,

    \begin{equation}
        \mathrm{Exactitud}= \frac{TP+TN}{TP+TN+FP+FN}.
        \label{eq:acc}
    \end{equation}


    \item Precisión: esta medida es la proporción de ejemplos pertenecientes a una clase que fueron correctamente predichos en dicha clase,
    
        \begin{equation}
            \mathrm{Precisi\acute{o}n} = \frac{TP}{TP+FP}.
            \label{eq:pre}
        \end{equation}
        
    \item Exhaustividad: esta medida es la proporción entre las predicciones que fueron correctamente clasificadas en una clase sobre los ejemplos que pertenecían a esa clase,
    
    \begin{equation}
        \mathrm{Exhaustividad}= \frac{TP}{TP+FN}.
        \label{eq:re}
    \end{equation}

    \item Métrica-F1: esta medida está diseñada para ponderar de manera equilibrada entre las métricas de precisión (\ref{eq:pre}) y exhaustividad (\ref{eq:re}),
    
    \begin{equation}
        \mathrm{F1} = \frac{2\cdot\mathrm{Precisi\acute{o}n}\cdot\mathrm{Exhaustividad}}{\mathrm{Precisi\acute{o}n}+\mathrm{Exhaustividad}}.
        \label{eq:f1}
    \end{equation}

\end{itemize}
